\section*{2. Objetivos}

\subsection*{2.1 Objetivo general}

El objetivo general de este Trabajo de Fin de Grado es el diseño e implementación de una 
plataforma web orientada al análisis predictivo y a la gestión de activos en el ámbito del 
Fantasy Football, con el fin de apoyar la toma de decisiones de los mánagers en un entorno 
caracterizado por la elevada complejidad y volumen de información.

Para ello, la plataforma se apoya en modelos de machine learning capaces de procesar datos 
históricos y actuales con el objetivo de generar predicciones sobre el rendimiento futuro de 
los jugadores y recomendaciones que contribuyan a optimizar la gestión de las plantillas, 
tanto en la selección de alineaciones como en la planificación de fichajes y ventas.

Un aspecto fundamental del objetivo general es la incorporación de técnicas de Explainable 
Artificial Intelligence (XAI), orientadas a dotar de transparencia a los modelos predictivos 
y a permitir que los usuarios comprendan los factores que influyen en cada recomendación, 
favoreciendo así la confianza y la validación de las decisiones propuestas por el sistema.

Asimismo, la plataforma se concibe con un enfoque de accesibilidad progresiva, combinando 
el acceso abierto a funcionalidades básicas sin necesidad de registro con un conjunto de 
funcionalidades avanzadas personalizadas para usuarios autenticados.

Finalmente, el sistema integra el análisis del contexto competitivo, incluyendo variables 
relacionadas con el próximo rival, con el objetivo de ofrecer predicciones y recomendaciones 
contextualizadas y alineadas con la realidad competitiva del Fantasy Football.

\subsection*{2.2 Objetivos}

Para la consecución del objetivo general, se definen los siguientes objetivos que describen 
las capacidades del sistema, y serán detallados más adelante en el apéndice IV 
(Requistios):

\begin{description}
\item[OBJ-F1:] Permitir la visualización de clasificación LaLiga y de los equipos de Fantasy 
Football sin necesidad de registro previo. Esta funcionalidad está orientada a ofrecer un 
acceso abierto a información relevante, facilitando que cualquier usuario pueda consultar 
datos básicos, estadísticas y análisis generales sin comprometer información personal. Esta 
decisión contribuye a mejorar la accesibilidad y a aumentar el alcance potencial de la 
plataforma.

\item[OBJ-F2:] Generación de predicciones sobre el rendimiento de los jugadores. Estas predicciones 
se basan en el análisis de datos históricos, estadísticas actuales y otros factores 
contextuales relevantes. El sistema debe ser capaz de transformar grandes volúmenes de 
información en estimaciones comprensibles que ayuden a los usuarios a anticipar el 
comportamiento futuro de los jugadores en las próximas jornadas.

\item[OBJ-F3:] Generar recomendaciones automáticas para la gestión de la plantilla de los usuarios 
registrados. Dichas recomendaciones abarcan aspectos clave del Fantasy Football, como la 
selección de alineaciones óptimas, la identificación de oportunidades de fichaje, la venta 
de jugadores con bajo rendimiento esperado o la sustitución estratégica de determinados 
activos. Este enfoque permite apoyar la toma de decisiones del usuario sin sustituir su 
criterio personal.

\item[OBJ-F4:] Análisis del próximo rival constituye otro objetivo funcional relevante del sistema. 
El rendimiento de un jugador no depende únicamente de sus estadísticas individuales, sino 
también del contexto del enfrentamiento. Por ello, la plataforma incorpora factores como la 
dificultad del partido, el estado de forma del rival o el rendimiento defensivo del equipo 
contrario, enriqueciendo así la calidad de las recomendaciones ofrecidas.

\item[OBJ-F5:] Para los usuarios autenticados, la plataforma ofrece funcionalidades avanzadas 
orientadas a la personalización de la experiencia. Entre ellas se incluye la gestión 
individualizada de equipos, permitiendo al usuario adaptar el sistema a su situación concreta 
dentro de la liga Fantasy.

\item[OBJ-F6:] Presentación de explicaciones claras y comprensibles sobre las recomendaciones 
generadas. El sistema no solo debe indicar qué acción se recomienda, sino también por qué se 
propone dicha acción, destacando los factores más influyentes en la decisión. Este enfoque 
refuerza la confianza del usuario y favorece un uso crítico de las recomendaciones 
automáticas.

\item[OBJ-F7:] Asegurar una interfaz web intuitiva y accesible, diseñada para ser utilizada tanto 
por usuarios con experiencia en Fantasy Football como por usuarios principiantes. La claridad 
visual, la organización de la información y la facilidad de navegación son aspectos clave para 
asegurar una experiencia de usuario satisfactoria.

\item[OBJ-F8:] Evaluar los modelos predictivos asegurando que superen un umbral mínimo de rendimiento 
según la métrica seleccionada (por ejemplo, MAE, RMSE o precisión). Esta evaluación garantiza 
que las predicciones sean confiables y útiles para la toma de decisiones, evitando incorporar 
modelos con resultados insuficientes o inconsistentes que puedan comprometer la validez de las 
recomendaciones. Se profundizará más en estos aspectos en el apéndice V 
(Implementación).

\item[OBJ-F9:] Ofrecer vistas detalladas de equipos y jugadores, incluyendo estadísticas históricas, 
rendimiento reciente, predicciones futuras y contexto competitivo, permitiendo al usuario 
profundizar en el análisis de cada activo.

\item[OBJ-F10:] Incorporar un apartado social que permita a los usuarios registrados añadir amigos, 
consultar sus plantillas y comparar estrategias, fomentando la interacción, la competitividad 
y el aprendizaje colaborativo entre mánagers.
\end{description}

\subsection*{2.3 Objetivos técnicos}

El objetivo técnico principal es la implementación de modelos de machine learning orientados a 
la predicción del rendimiento de los jugadores y a la generación de recomendaciones. Estos 
modelos deben ser capaces de procesar datos deportivos reales procedentes de distintas fuentes, 
adaptarse a la evolución temporal de la competición y ofrecer resultados consistentes y 
actualizados que sirvan de apoyo a la toma de decisiones de los usuarios.

La obtención de los datos necesarios para alimentar estos modelos constituye en sí misma un 
objetivo técnico relevante. Para ello, el sistema debe incorporar mecanismos de web scraping 
sobre fuentes deportivas fiables y contrastadas, permitiendo la recopilación automática de 
estadísticas, información de jugadores y contexto competitivo. Este proceso conlleva una 
complejidad adicional, ya que cada fuente presenta su propia estructura, formato y criterios 
de nomenclatura.

En particular, uno de los principales retos técnicos asociados al uso de múltiples fuentes de 
información es la unificación de los datos relativos a los jugadores. Diferentes sitios web 
pueden emplear nombres distintos para un mismo jugador, abreviaturas, variaciones ortográficas 
o incluso identificadores propios. Por este motivo, el proyecto contempla el diseño e 
implementación de un mecanismo de matching y normalización de entidades que permita identificar 
de forma inequívoca a cada jugador y poner en común la información procedente de diversas 
fuentes. Este proceso resulta fundamental para garantizar la coherencia de los datos y evitar 
inconsistencias que puedan afectar al rendimiento de los modelos predictivos.

La integración de técnicas de Explainable Artificial Intelligence (XAI) constituye otro 
objetivo técnico clave del proyecto. Se busca incorporar métodos que permitan analizar la 
relevancia de las características utilizadas por los modelos y ofrecer explicaciones 
interpretables para los usuarios finales. Esta integración supone un reto adicional, ya que 
combina aspectos avanzados de inteligencia artificial con requisitos de usabilidad y 
comunicación de resultados, además de tratarse de una tecnología que no había sido utilizada 
previamente.

Asimismo, se plantea como objetivo la implementación de un sistema de autenticación y 
autorización que permita diferenciar entre usuarios anónimos y usuarios registrados, 
garantizando un acceso controlado a las funcionalidades avanzadas y protegiendo la información 
personal de los usuarios.

La gestión y el procesamiento de datos deportivos reales representan otro objetivo técnico 
relevante. El sistema debe asegurar la correcta limpieza, consistencia y actualización de los 
datos, ya que la calidad de la información de entrada influye directamente en la fiabilidad de 
las predicciones generadas y en la validez de las recomendaciones ofrecidas.

Se establece como objetivo técnico la validación y evaluación del sistema desarrollado, 
definiendo métricas adecuadas que permitan medir la calidad de las predicciones y analizar el 
comportamiento de los modelos en distintos escenarios. Este proceso de evaluación resulta 
esencial para justificar las decisiones técnicas adoptadas y para identificar las fortalezas y 
limitaciones del sistema.

Por último, se establecen como objetivos técnicos la optimización del rendimiento del sistema y 
la garantía de la seguridad de los datos. Esto incluye la reducción de tiempos de respuesta, la 
gestión eficiente de recursos y la protección de la información almacenada, aspectos 
fundamentales garantizando la viabilidad y la calidad del sistema desarrollado, asumiendo la 
existencia de incertidumbre inherente a los datos deportivos y a los eventos impredecibles 
propios del dominio.
